% Circuit: 2nd-order Strang splitting Trotter step for 3-qubit periodic Heisenberg chain
% Generated for: QuantumFurnace.jl PhD thesis
% Usage: % Circuit: 2nd-order Strang splitting Trotter step for 3-qubit periodic Heisenberg chain
% Generated for: QuantumFurnace.jl PhD thesis
% Usage: % Circuit: 2nd-order Strang splitting Trotter step for 3-qubit periodic Heisenberg chain
% Generated for: QuantumFurnace.jl PhD thesis
% Usage: % Circuit: 2nd-order Strang splitting Trotter step for 3-qubit periodic Heisenberg chain
% Generated for: QuantumFurnace.jl PhD thesis
% Usage: \input{circuits/trotter-strang-3q.tex} in main thesis, or compile standalone
%
% Matches src/trotter_domain.jl:_trotterize2 for n=3 periodic.
% Bond groups: odd (1,2), even (2,3), boundary (3,1).
% Each bond Hamiltonian H_{jk} = J_x X_jX_k + J_y Y_jY_k + J_z Z_jZ_k is further
% decomposed into three commuting two-qubit exponentials. Because [X_jX_k, Y_jY_k]
% = [Y_jY_k, Z_jZ_k] = [Z_jZ_k, X_jX_k] = 0 on the same bond, this factorisation
% is exact — it adds no additional Trotter error, and each factor is a native
% two-qubit gate on most hardware (R_{XX}, R_{YY}, R_{ZZ}).
%
% Strang palindrome (combining adjacent boundary terms):
%   e^{-iH_{12} dt/2} e^{-iH_{23} dt/2} e^{-iH_{31} dt} e^{-iH_{23} dt/2} e^{-iH_{12} dt/2}
%
\documentclass{article}
\usepackage[margin=1cm]{geometry}
\usepackage{tikz}
\usetikzlibrary{quantikz2}
\usepackage{amsmath,amssymb}
\usepackage{braket}
\usepackage{graphicx}

\newcommand{\ee}{\mathrm{e}}
\newcommand{\ii}{\mathrm{i}}

% Rounded corners on all gates globally
\tikzset{operator/.append style={rounded corners}}

% NOTE: \dashedthrough (and its `dashed line` tikz style) is assumed to be
% defined already in the including thesis preamble. Re-declare here only
% if you want this file to compile standalone.

\begin{document}
\pagestyle{empty}

\begin{figure}[ht]
  \centering
  % [transparent] makes gate boxes see-through so \dashedthrough is visible
  % inside the 3-qubit boundary gates. \resizebox ensures the figure always
  % fits \textwidth regardless of the (inherently wide) 15-column layout.
  % Tight gate padding is applied *locally* via /tikz/operator/.append style,
  % so it does not affect other circuits in the thesis.
  % Local tweaks (do not leak to other circuits):
  %  * font=\scriptsize inside gate boxes shrinks natural box width ~30%
  %  * column sep is large enough that gates don't overlap at natural size
  %    (resizebox preserves overlap, so non-overlap must hold pre-scale)
  %  * resizebox then scales the whole picture down to \textwidth
  % Label size is tuned relative to natural picture width:
  %  * \resizebox scales the whole picture uniformly, so growing font AND
  %    column sep together does nothing. To get bigger labels in the final
  %    PDF, grow the font more than the column spacing.
  %  * \small (9pt) renders noticeably larger than \scriptsize (7pt); we
  %    keep column sep modest so \resizebox does not over-shrink.
  \resizebox{\textwidth}{!}{%
  \begin{quantikz}[
    transparent, row sep=0.6cm, column sep=0.4cm,
    /tikz/operator/.append style={inner sep=3pt, minimum height=6mm, font=\small}
  ]
  %
  % Column layout (15 gate cols + final \qw = 16 cells per row):
  %   cols  1-3 : A-block L  = H_{12} dt/2  (q_1, q_2)
  %   cols  4-6 : B-block L  = H_{23} dt/2  (q_2, q_3)
  %   cols  7-9 : C-block    = H_{31} dt    (q_1, q_3; q_2 dashed)
  %   cols 10-12: B-block R  = H_{23} dt/2  (q_2, q_3)
  %   cols 13-15: A-block R  = H_{12} dt/2  (q_1, q_2)
  % All rows must have the same number of cells; inactive bond groups
  % on a given row are filled with plain \qw (empty cells under a
  % \gate[n]{...} are consumed by the multi-row gate box).
  %
    \lstick{$q_1$}
  % A-block L (active)
      & \gate[2]{R_{X_1X_2}}
      & \gate[2]{R_{Y_1Y_2}}
      & \gate[2]{R_{Z_1Z_2}}
  % B-block L (q_1 idle)
      & \qw & \qw & \qw
  % C-block (boundary, q_1 leads the [3]-wide gate; R^2 = R applied twice = rotate by 2 dt)
      & \gate[3, label style={yshift=0.3cm}]{R_{X_3X_1}^{2}}
      & \gate[3, label style={yshift=0.3cm}]{R_{Y_3Y_1}^{2}}
      & \gate[3, label style={yshift=0.3cm}]{R_{Z_3Z_1}^{2}}
  % B-block R (q_1 idle)
      & \qw & \qw & \qw
  % A-block R (active)
      & \gate[2]{R_{X_1X_2}}
      & \gate[2]{R_{Y_1Y_2}}
      & \gate[2]{R_{Z_1Z_2}}
      & \qw \\
  %
    \lstick{$q_2$}
  % A-block L (empty: consumed by q_1's gate[2])
      & & &
  % B-block L (active, q_2 leads)
      & \gate[2]{R_{X_2X_3}}
      & \gate[2]{R_{Y_2Y_3}}
      & \gate[2]{R_{Z_2Z_3}}
  % C-block: q_2 passes through unaffected
      & \dashedthrough & \dashedthrough & \dashedthrough
  % B-block R (active, q_2 leads)
      & \gate[2]{R_{X_2X_3}}
      & \gate[2]{R_{Y_2Y_3}}
      & \gate[2]{R_{Z_2Z_3}}
  % A-block R (empty: consumed by q_1's gate[2])
      & & &
      & \qw \\
  %
    \lstick{$q_3$}
  % A-block L (q_3 idle)
      & \qw & \qw & \qw
  % B-block L (empty: consumed by q_2's gate[2])
      & & &
  % C-block (empty: consumed by q_1's gate[3])
      & & &
  % B-block R (empty: consumed by q_2's gate[2])
      & & &
  % A-block R (q_3 idle)
      & \qw & \qw & \qw
      & \qw
  \end{quantikz}%
  }
  \caption{One second-order Strang Trotter step for the $n{=}3$ periodic
  isotropic Heisenberg chain with $\delta t = t/r$, using the hardware notation
  $R_{\sigma_j\sigma_k}(\theta) \equiv \ee^{-\ii\theta\,\sigma_j\sigma_k/2}$ and
  the abbreviation $R_{\sigma_j\sigma_k} \equiv R_{\sigma_j\sigma_k}(\delta t)$.
  The boundary bond $(3,1)$ closes the ring and $q_2$ passes through unaffected
  (dashed wire).}
  \label{fig:trotter-strang-3q}
\end{figure}

\end{document}
 in main thesis, or compile standalone
%
% Matches src/trotter_domain.jl:_trotterize2 for n=3 periodic.
% Bond groups: odd (1,2), even (2,3), boundary (3,1).
% Each bond Hamiltonian H_{jk} = J_x X_jX_k + J_y Y_jY_k + J_z Z_jZ_k is further
% decomposed into three commuting two-qubit exponentials. Because [X_jX_k, Y_jY_k]
% = [Y_jY_k, Z_jZ_k] = [Z_jZ_k, X_jX_k] = 0 on the same bond, this factorisation
% is exact — it adds no additional Trotter error, and each factor is a native
% two-qubit gate on most hardware (R_{XX}, R_{YY}, R_{ZZ}).
%
% Strang palindrome (combining adjacent boundary terms):
%   e^{-iH_{12} dt/2} e^{-iH_{23} dt/2} e^{-iH_{31} dt} e^{-iH_{23} dt/2} e^{-iH_{12} dt/2}
%
\documentclass{article}
\usepackage[margin=1cm]{geometry}
\usepackage{tikz}
\usetikzlibrary{quantikz2}
\usepackage{amsmath,amssymb}
\usepackage{braket}
\usepackage{graphicx}

\newcommand{\ee}{\mathrm{e}}
\newcommand{\ii}{\mathrm{i}}

% Rounded corners on all gates globally
\tikzset{operator/.append style={rounded corners}}

% NOTE: \dashedthrough (and its `dashed line` tikz style) is assumed to be
% defined already in the including thesis preamble. Re-declare here only
% if you want this file to compile standalone.

\begin{document}
\pagestyle{empty}

\begin{figure}[ht]
  \centering
  % [transparent] makes gate boxes see-through so \dashedthrough is visible
  % inside the 3-qubit boundary gates. \resizebox ensures the figure always
  % fits \textwidth regardless of the (inherently wide) 15-column layout.
  % Tight gate padding is applied *locally* via /tikz/operator/.append style,
  % so it does not affect other circuits in the thesis.
  % Local tweaks (do not leak to other circuits):
  %  * font=\scriptsize inside gate boxes shrinks natural box width ~30%
  %  * column sep is large enough that gates don't overlap at natural size
  %    (resizebox preserves overlap, so non-overlap must hold pre-scale)
  %  * resizebox then scales the whole picture down to \textwidth
  % Label size is tuned relative to natural picture width:
  %  * \resizebox scales the whole picture uniformly, so growing font AND
  %    column sep together does nothing. To get bigger labels in the final
  %    PDF, grow the font more than the column spacing.
  %  * \small (9pt) renders noticeably larger than \scriptsize (7pt); we
  %    keep column sep modest so \resizebox does not over-shrink.
  \resizebox{\textwidth}{!}{%
  \begin{quantikz}[
    transparent, row sep=0.6cm, column sep=0.4cm,
    /tikz/operator/.append style={inner sep=3pt, minimum height=6mm, font=\small}
  ]
  %
  % Column layout (15 gate cols + final \qw = 16 cells per row):
  %   cols  1-3 : A-block L  = H_{12} dt/2  (q_1, q_2)
  %   cols  4-6 : B-block L  = H_{23} dt/2  (q_2, q_3)
  %   cols  7-9 : C-block    = H_{31} dt    (q_1, q_3; q_2 dashed)
  %   cols 10-12: B-block R  = H_{23} dt/2  (q_2, q_3)
  %   cols 13-15: A-block R  = H_{12} dt/2  (q_1, q_2)
  % All rows must have the same number of cells; inactive bond groups
  % on a given row are filled with plain \qw (empty cells under a
  % \gate[n]{...} are consumed by the multi-row gate box).
  %
    \lstick{$q_1$}
  % A-block L (active)
      & \gate[2]{R_{X_1X_2}}
      & \gate[2]{R_{Y_1Y_2}}
      & \gate[2]{R_{Z_1Z_2}}
  % B-block L (q_1 idle)
      & \qw & \qw & \qw
  % C-block (boundary, q_1 leads the [3]-wide gate; R^2 = R applied twice = rotate by 2 dt)
      & \gate[3, label style={yshift=0.3cm}]{R_{X_3X_1}^{2}}
      & \gate[3, label style={yshift=0.3cm}]{R_{Y_3Y_1}^{2}}
      & \gate[3, label style={yshift=0.3cm}]{R_{Z_3Z_1}^{2}}
  % B-block R (q_1 idle)
      & \qw & \qw & \qw
  % A-block R (active)
      & \gate[2]{R_{X_1X_2}}
      & \gate[2]{R_{Y_1Y_2}}
      & \gate[2]{R_{Z_1Z_2}}
      & \qw \\
  %
    \lstick{$q_2$}
  % A-block L (empty: consumed by q_1's gate[2])
      & & &
  % B-block L (active, q_2 leads)
      & \gate[2]{R_{X_2X_3}}
      & \gate[2]{R_{Y_2Y_3}}
      & \gate[2]{R_{Z_2Z_3}}
  % C-block: q_2 passes through unaffected
      & \dashedthrough & \dashedthrough & \dashedthrough
  % B-block R (active, q_2 leads)
      & \gate[2]{R_{X_2X_3}}
      & \gate[2]{R_{Y_2Y_3}}
      & \gate[2]{R_{Z_2Z_3}}
  % A-block R (empty: consumed by q_1's gate[2])
      & & &
      & \qw \\
  %
    \lstick{$q_3$}
  % A-block L (q_3 idle)
      & \qw & \qw & \qw
  % B-block L (empty: consumed by q_2's gate[2])
      & & &
  % C-block (empty: consumed by q_1's gate[3])
      & & &
  % B-block R (empty: consumed by q_2's gate[2])
      & & &
  % A-block R (q_3 idle)
      & \qw & \qw & \qw
      & \qw
  \end{quantikz}%
  }
  \caption{One second-order Strang Trotter step for the $n{=}3$ periodic
  isotropic Heisenberg chain with $\delta t = t/r$, using the hardware notation
  $R_{\sigma_j\sigma_k}(\theta) \equiv \ee^{-\ii\theta\,\sigma_j\sigma_k/2}$ and
  the abbreviation $R_{\sigma_j\sigma_k} \equiv R_{\sigma_j\sigma_k}(\delta t)$.
  The boundary bond $(3,1)$ closes the ring and $q_2$ passes through unaffected
  (dashed wire).}
  \label{fig:trotter-strang-3q}
\end{figure}

\end{document}
 in main thesis, or compile standalone
%
% Matches src/trotter_domain.jl:_trotterize2 for n=3 periodic.
% Bond groups: odd (1,2), even (2,3), boundary (3,1).
% Each bond Hamiltonian H_{jk} = J_x X_jX_k + J_y Y_jY_k + J_z Z_jZ_k is further
% decomposed into three commuting two-qubit exponentials. Because [X_jX_k, Y_jY_k]
% = [Y_jY_k, Z_jZ_k] = [Z_jZ_k, X_jX_k] = 0 on the same bond, this factorisation
% is exact — it adds no additional Trotter error, and each factor is a native
% two-qubit gate on most hardware (R_{XX}, R_{YY}, R_{ZZ}).
%
% Strang palindrome (combining adjacent boundary terms):
%   e^{-iH_{12} dt/2} e^{-iH_{23} dt/2} e^{-iH_{31} dt} e^{-iH_{23} dt/2} e^{-iH_{12} dt/2}
%
\documentclass{article}
\usepackage[margin=1cm]{geometry}
\usepackage{tikz}
\usetikzlibrary{quantikz2}
\usepackage{amsmath,amssymb}
\usepackage{braket}
\usepackage{graphicx}

\newcommand{\ee}{\mathrm{e}}
\newcommand{\ii}{\mathrm{i}}

% Rounded corners on all gates globally
\tikzset{operator/.append style={rounded corners}}

% NOTE: \dashedthrough (and its `dashed line` tikz style) is assumed to be
% defined already in the including thesis preamble. Re-declare here only
% if you want this file to compile standalone.

\begin{document}
\pagestyle{empty}

\begin{figure}[ht]
  \centering
  % [transparent] makes gate boxes see-through so \dashedthrough is visible
  % inside the 3-qubit boundary gates. \resizebox ensures the figure always
  % fits \textwidth regardless of the (inherently wide) 15-column layout.
  % Tight gate padding is applied *locally* via /tikz/operator/.append style,
  % so it does not affect other circuits in the thesis.
  % Local tweaks (do not leak to other circuits):
  %  * font=\scriptsize inside gate boxes shrinks natural box width ~30%
  %  * column sep is large enough that gates don't overlap at natural size
  %    (resizebox preserves overlap, so non-overlap must hold pre-scale)
  %  * resizebox then scales the whole picture down to \textwidth
  % Label size is tuned relative to natural picture width:
  %  * \resizebox scales the whole picture uniformly, so growing font AND
  %    column sep together does nothing. To get bigger labels in the final
  %    PDF, grow the font more than the column spacing.
  %  * \small (9pt) renders noticeably larger than \scriptsize (7pt); we
  %    keep column sep modest so \resizebox does not over-shrink.
  \resizebox{\textwidth}{!}{%
  \begin{quantikz}[
    transparent, row sep=0.6cm, column sep=0.4cm,
    /tikz/operator/.append style={inner sep=3pt, minimum height=6mm, font=\small}
  ]
  %
  % Column layout (15 gate cols + final \qw = 16 cells per row):
  %   cols  1-3 : A-block L  = H_{12} dt/2  (q_1, q_2)
  %   cols  4-6 : B-block L  = H_{23} dt/2  (q_2, q_3)
  %   cols  7-9 : C-block    = H_{31} dt    (q_1, q_3; q_2 dashed)
  %   cols 10-12: B-block R  = H_{23} dt/2  (q_2, q_3)
  %   cols 13-15: A-block R  = H_{12} dt/2  (q_1, q_2)
  % All rows must have the same number of cells; inactive bond groups
  % on a given row are filled with plain \qw (empty cells under a
  % \gate[n]{...} are consumed by the multi-row gate box).
  %
    \lstick{$q_1$}
  % A-block L (active)
      & \gate[2]{R_{X_1X_2}}
      & \gate[2]{R_{Y_1Y_2}}
      & \gate[2]{R_{Z_1Z_2}}
  % B-block L (q_1 idle)
      & \qw & \qw & \qw
  % C-block (boundary, q_1 leads the [3]-wide gate; R^2 = R applied twice = rotate by 2 dt)
      & \gate[3, label style={yshift=0.3cm}]{R_{X_3X_1}^{2}}
      & \gate[3, label style={yshift=0.3cm}]{R_{Y_3Y_1}^{2}}
      & \gate[3, label style={yshift=0.3cm}]{R_{Z_3Z_1}^{2}}
  % B-block R (q_1 idle)
      & \qw & \qw & \qw
  % A-block R (active)
      & \gate[2]{R_{X_1X_2}}
      & \gate[2]{R_{Y_1Y_2}}
      & \gate[2]{R_{Z_1Z_2}}
      & \qw \\
  %
    \lstick{$q_2$}
  % A-block L (empty: consumed by q_1's gate[2])
      & & &
  % B-block L (active, q_2 leads)
      & \gate[2]{R_{X_2X_3}}
      & \gate[2]{R_{Y_2Y_3}}
      & \gate[2]{R_{Z_2Z_3}}
  % C-block: q_2 passes through unaffected
      & \dashedthrough & \dashedthrough & \dashedthrough
  % B-block R (active, q_2 leads)
      & \gate[2]{R_{X_2X_3}}
      & \gate[2]{R_{Y_2Y_3}}
      & \gate[2]{R_{Z_2Z_3}}
  % A-block R (empty: consumed by q_1's gate[2])
      & & &
      & \qw \\
  %
    \lstick{$q_3$}
  % A-block L (q_3 idle)
      & \qw & \qw & \qw
  % B-block L (empty: consumed by q_2's gate[2])
      & & &
  % C-block (empty: consumed by q_1's gate[3])
      & & &
  % B-block R (empty: consumed by q_2's gate[2])
      & & &
  % A-block R (q_3 idle)
      & \qw & \qw & \qw
      & \qw
  \end{quantikz}%
  }
  \caption{One second-order Strang Trotter step for the $n{=}3$ periodic
  isotropic Heisenberg chain with $\delta t = t/r$, using the hardware notation
  $R_{\sigma_j\sigma_k}(\theta) \equiv \ee^{-\ii\theta\,\sigma_j\sigma_k/2}$ and
  the abbreviation $R_{\sigma_j\sigma_k} \equiv R_{\sigma_j\sigma_k}(\delta t)$.
  The boundary bond $(3,1)$ closes the ring and $q_2$ passes through unaffected
  (dashed wire).}
  \label{fig:trotter-strang-3q}
\end{figure}

\end{document}
 in main thesis, or compile standalone
%
% Matches src/trotter_domain.jl:_trotterize2 for n=3 periodic.
% Bond groups: odd (1,2), even (2,3), boundary (3,1).
% Each bond Hamiltonian H_{jk} = J_x X_jX_k + J_y Y_jY_k + J_z Z_jZ_k is further
% decomposed into three commuting two-qubit exponentials. Because [X_jX_k, Y_jY_k]
% = [Y_jY_k, Z_jZ_k] = [Z_jZ_k, X_jX_k] = 0 on the same bond, this factorisation
% is exact — it adds no additional Trotter error, and each factor is a native
% two-qubit gate on most hardware (R_{XX}, R_{YY}, R_{ZZ}).
%
% Strang palindrome (combining adjacent boundary terms):
%   e^{-iH_{12} dt/2} e^{-iH_{23} dt/2} e^{-iH_{31} dt} e^{-iH_{23} dt/2} e^{-iH_{12} dt/2}
%
\documentclass{article}
\usepackage[margin=1cm]{geometry}
\usepackage{tikz}
\usetikzlibrary{quantikz2}
\usepackage{amsmath,amssymb}
\usepackage{braket}
\usepackage{graphicx}

\newcommand{\ee}{\mathrm{e}}
\newcommand{\ii}{\mathrm{i}}

% Rounded corners on all gates globally
\tikzset{operator/.append style={rounded corners}}

% NOTE: \dashedthrough (and its `dashed line` tikz style) is assumed to be
% defined already in the including thesis preamble. Re-declare here only
% if you want this file to compile standalone.

\begin{document}
\pagestyle{empty}

\begin{figure}[ht]
  \centering
  % [transparent] makes gate boxes see-through so \dashedthrough is visible
  % inside the 3-qubit boundary gates. \resizebox ensures the figure always
  % fits \textwidth regardless of the (inherently wide) 15-column layout.
  % Tight gate padding is applied *locally* via /tikz/operator/.append style,
  % so it does not affect other circuits in the thesis.
  % Local tweaks (do not leak to other circuits):
  %  * font=\scriptsize inside gate boxes shrinks natural box width ~30%
  %  * column sep is large enough that gates don't overlap at natural size
  %    (resizebox preserves overlap, so non-overlap must hold pre-scale)
  %  * resizebox then scales the whole picture down to \textwidth
  % Label size is tuned relative to natural picture width:
  %  * \resizebox scales the whole picture uniformly, so growing font AND
  %    column sep together does nothing. To get bigger labels in the final
  %    PDF, grow the font more than the column spacing.
  %  * \small (9pt) renders noticeably larger than \scriptsize (7pt); we
  %    keep column sep modest so \resizebox does not over-shrink.
  \resizebox{\textwidth}{!}{%
  \begin{quantikz}[
    transparent, row sep=0.6cm, column sep=0.4cm,
    /tikz/operator/.append style={inner sep=3pt, minimum height=6mm, font=\small}
  ]
  %
  % Column layout (15 gate cols + final \qw = 16 cells per row):
  %   cols  1-3 : A-block L  = H_{12} dt/2  (q_1, q_2)
  %   cols  4-6 : B-block L  = H_{23} dt/2  (q_2, q_3)
  %   cols  7-9 : C-block    = H_{31} dt    (q_1, q_3; q_2 dashed)
  %   cols 10-12: B-block R  = H_{23} dt/2  (q_2, q_3)
  %   cols 13-15: A-block R  = H_{12} dt/2  (q_1, q_2)
  % All rows must have the same number of cells; inactive bond groups
  % on a given row are filled with plain \qw (empty cells under a
  % \gate[n]{...} are consumed by the multi-row gate box).
  %
    \lstick{$q_1$}
  % A-block L (active)
      & \gate[2]{R_{X_1X_2}}
      & \gate[2]{R_{Y_1Y_2}}
      & \gate[2]{R_{Z_1Z_2}}
  % B-block L (q_1 idle)
      & \qw & \qw & \qw
  % C-block (boundary, q_1 leads the [3]-wide gate; R^2 = R applied twice = rotate by 2 dt)
      & \gate[3, label style={yshift=0.3cm}]{R_{X_3X_1}^{2}}
      & \gate[3, label style={yshift=0.3cm}]{R_{Y_3Y_1}^{2}}
      & \gate[3, label style={yshift=0.3cm}]{R_{Z_3Z_1}^{2}}
  % B-block R (q_1 idle)
      & \qw & \qw & \qw
  % A-block R (active)
      & \gate[2]{R_{X_1X_2}}
      & \gate[2]{R_{Y_1Y_2}}
      & \gate[2]{R_{Z_1Z_2}}
      & \qw \\
  %
    \lstick{$q_2$}
  % A-block L (empty: consumed by q_1's gate[2])
      & & &
  % B-block L (active, q_2 leads)
      & \gate[2]{R_{X_2X_3}}
      & \gate[2]{R_{Y_2Y_3}}
      & \gate[2]{R_{Z_2Z_3}}
  % C-block: q_2 passes through unaffected
      & \dashedthrough & \dashedthrough & \dashedthrough
  % B-block R (active, q_2 leads)
      & \gate[2]{R_{X_2X_3}}
      & \gate[2]{R_{Y_2Y_3}}
      & \gate[2]{R_{Z_2Z_3}}
  % A-block R (empty: consumed by q_1's gate[2])
      & & &
      & \qw \\
  %
    \lstick{$q_3$}
  % A-block L (q_3 idle)
      & \qw & \qw & \qw
  % B-block L (empty: consumed by q_2's gate[2])
      & & &
  % C-block (empty: consumed by q_1's gate[3])
      & & &
  % B-block R (empty: consumed by q_2's gate[2])
      & & &
  % A-block R (q_3 idle)
      & \qw & \qw & \qw
      & \qw
  \end{quantikz}%
  }
  \caption{One second-order Strang Trotter step for the $n{=}3$ periodic
  isotropic Heisenberg chain with $\delta t = t/r$, using the hardware notation
  $R_{\sigma_j\sigma_k}(\theta) \equiv \ee^{-\ii\theta\,\sigma_j\sigma_k/2}$ and
  the abbreviation $R_{\sigma_j\sigma_k} \equiv R_{\sigma_j\sigma_k}(\delta t)$.
  The boundary bond $(3,1)$ closes the ring and $q_2$ passes through unaffected
  (dashed wire).}
  \label{fig:trotter-strang-3q}
\end{figure}

\end{document}
